% !TEX program = xelatex
\documentclass{chinese-erj}

% ── 额外宏包 ──
\usepackage{booktabs}
\usepackage{threeparttable}
\usepackage{tabularx}
\usepackage{geometry}
\usepackage{graphicx}
\usepackage{subfigure}
\usepackage{appendix}
\usepackage{hyperref}
\usepackage{amsmath}
\usepackage{amssymb}
\usepackage{bm}
\usepackage{multirow}
\usepackage{array}
\usepackage{longtable}
\usepackage{makecell}
\usepackage{pifont}
\usepackage{caption}
\usepackage{float}
\usepackage{color}
\usepackage{xcolor}
\usepackage{enumitem}
\setlist{nosep, left=2em}


% ── 论文信息 ──
\title{信息披露质量、多源特征融合与监管问询概率预测\thanks{本文基于"东吴证券杯"第五届中国研究生金融科技创新大赛数据集完成。感谢东吴证券股份有限公司提供的数据与业务场景支持。文责自负。}}

\author{作者姓名}
\date{}

\header{信息披露质量、多源特征融合与监管问询概率预测}

% 英文信息
\entitle{Information Disclosure Quality, Multi-source Feature Fusion, and Regulatory Inquiry Probability Prediction}
\enauthor{Author Name}
\eninstitute{School of Economics, University Name}
\enabstract{
This paper constructs a listed company regulatory inquiry probability prediction model by integrating multi-source heterogeneous information, including financial indicators, LLM-extracted textual semantic features, market microstructure characteristics, and corporate governance variables. Based on a dataset of approximately 100,000 firm-quarter observations from Chinese A-share listed companies, we systematically examine the incremental predictive power of "soft information" (textual features) relative to traditional "hard information" (financial indicators). Our logistic regression analysis reveals that traditional financial anomaly indicators possess significant predictive power, while large language model-extracted textual semantic features provide substantial incremental predictive value, increasing the AUC from 0.71 to 0.77 -- a statistically significant improvement confirmed by the DeLong test (p<0.01). The economic significance of this increment is most pronounced among firms with severe asymmetric information, including small-cap firms and those with limited analyst coverage. Further using LightGBM heterogeneous stacking ensemble, the full model achieves an AUC of 0.79 and a Top-10\% coverage of 38.6\%. Combining SHAP-based interpretable attribution analysis with large language model-generated explanatory reports, we achieve a three-level traceable attribution chain from probabilistic output to raw evidence, and the explanation quality score reaches 86.3 points measured by LLM-as-Judge. This research provides empirical evidence for understanding the behavioral logic of regulatory inquiries and offers directional references for intelligent financial supervision tools.
}
\enkeywords{Regulatory Inquiry; Feature Fusion; Large Language Model; Explainable AI; Information Asymmetry}
\jel{C53, G14, G18, M41}


\begin{document}

\maketitle

\begin{abstract}
上市公司监管问询是交易所事中事后监管的核心工具，但哪些公司更可能被问询、不同信息维度的预测价值如何、预测结果能否提供可解释的经济学归因，仍缺乏系统性研究。本文融合财务指标、大语言模型（LLM）提取的文本语义特征、市场微观结构特征和公司治理变量等多源异构信息，构建了面向中国A股上市公司的监管问询概率预测模型，系统检验了"软信息"相对于传统"硬信息"的增量预测价值。基于约10万公司-季度观测数据的实证研究发现：第一，传统财务异常指标（Beneish M-Score、应收增速偏离等）对问询概率具有显著的预测能力，基准Logistic回归AUC达到0.71；第二，LLM提取的文本语义风险特征提供了显著的增量信息价值，将AUC进一步提升至0.77，Delong检验表明AUC提升在1\%水平上统计显著；第三，软信息的增量价值在信息不对称更为严重的小市值公司和低分析师覆盖公司中更为凸显；第四，LightGBM异质集成全模型AUC达0.79，Top-10\%高风险公司覆盖率达38.6\%；第五，结合SHAP可解释归因分析与LLM生成式解释报告，实现了从概率输出到原文证据的三级可追溯归因链，LLM-as-Judge评估的逻辑解释有效性得分为86.3分。本文为理解监管问询的行为逻辑提供了经济学证据，也为智能化金融监管工具的开发提供了方向性参考。
\keywords{监管问询；多源特征融合；大语言模型；可解释人工智能；信息不对称}
\end{abstract}


% ── 引入各章节 ──
% ═══════════════════════════════════════════════════════════════
% 第一章：引言
% ═══════════════════════════════════════════════════════════════

\section{引 言}

监管问询是中国资本市场信息披露制度中事中事后监管的核心工具。自2014年上海证券交易所和深圳证券交易所全面推行问询函制度以来，问询函的发放数量和覆盖范围持续扩大。2023年全年，两大交易所发出各类问询函超过1500份，覆盖超过三分之一的A股上市公司。对于上市公司而言，收到问询函不仅意味着信息披露瑕疵被公开曝光，还可能引发股价下跌\citep{chen2019supervision}、分析师评级下调\citep{wang2019analyst}以及投资者诉讼等一系列负面后果。因此，在问询函正式发出之前提前预判哪些公司可能"被问询"，对于投资风险管理和监管资源配置均具有重要的实践意义。

从学术研究的角度看，监管问询的可预测性问题蕴涵着重要的经济学含义。监管者的问询行为并非随机——它反映了监管者在信息不对称环境下对上市公司信息披露质量的系统性判断。已有文献从财务指标异常\citep{cassell2013review}、公司治理缺陷\citep{chen2019effective}和信息披露质量\citep{li2019inquiry}等角度考察了问询概率的驱动因素，但受限于数据和方法，现有研究存在三个方面的明显不足。

第一个不足是数据维度单一。现有预测研究几乎完全依赖结构化财务指标（如Beneish M-Score、资产负债率等），忽略了上市公司公告文本中蕴含的丰富语义信息。事实上，公告文本中的管理层语调、信息披露充分性、MD\&A前瞻性陈述等"软信息"维度，可能包含财务指标无法反映的重要风险信号。正如Liberti和Petersen（2019）在《经济文献杂志》的综述中所强调的，软信息和硬信息在风险评估中具有互补性，二者的融合能够带来比单一信息维度更优的预测表现\citep{liberti2019information}。

第二个不足是工具手段有限。大语言模型（Large Language Model, LLM）在语义理解方面的突破性进展为从非结构化文本中提取结构化信息提供了全新的技术手段\citep{huang2023finbert}。然而，LLM在监管问询预测中的应用尚属空白——现有文献既未系统检验LLM提取的语义特征对问询概率的增量预测能力，也未比较其相对于传统财务指标的相对重要性。

第三个不足是可解释性缺失。金融监管场景对AI模型的可解释性和可审计性有极高要求——投研和风控人员不仅需要知道"概率是多少"，更需要理解"为什么是这个概率"。现有预测模型多为"黑箱"式输出，缺乏从概率分数到风险特征再到原文证据的完整归因链。SHAP（Shapley Additive Explanations）等可解释性方法在金融领域的计量经济学应用中具有独特价值\citep{lu2023ml}，但其与大型语言模型生成式归因报告的融合在监管问询场景中尚缺乏系统探索。

本文在上述三个维度上做出贡献。在数据维度，本文构建了涵盖六大类约235个特征的多元信息融合框架（财务特征、LLM语义特征、市场特征、历史监管特征、知识图谱特征、时序衍生特征），将结构化与非结构化信息纳入统一分析框架。在方法维度，本文首次将大语言模型提取的文本语义风险特征引入监管问询概率预测，并采用Logistic回归、嵌套模型比较（Nested Model Comparison）\citep{delong1988comparing}和异构集成学习相结合的策略，系统检验软信息的增量预测价值及其在信息不对称不同公司间的异质性。在可解释性维度，本文构建了"SHAP数值归因—LLM生成式归因报告—原文证据追溯"三级可解释机制，并通过LLM-as-Judge框架对归因质量进行多维度量化评估\citep{zheng2024judge}。

基于中国A股上市公司约10万条公司-季度观测数据的实证分析表明：第一，传统财务异常指标对问询概率具有显著的预测能力（Logistic AUC 0.71）；第二，LLM提取的文本语义特征提供了显著且经济上重要的增量信息（全模型AUC 0.78，Delong检验p$<$0.01）；第三，软信息的增量价值在信息不对称更严重的小市值公司和低分析师覆盖公司中更加凸显；第四，LightGBM异质集成全模型的AUC达到0.79，Top-10\%高风险公司的覆盖率达38.6\%；第五，结合SHAP和LLM的三级可解释报告在LLM-as-Judge评估中获得86.3分的综合质量评价。

本文的研究发现对于完善事中事后监管制度、提升监管效率具有明确的政策含义：监管机构可参考本文融合多源信息的预测框架，将文本语义分析纳入日常监管工具箱，构建"财务指标+文本语义+市场信号"三位一体的智能化问询筛选系统。同时，可解释的归因报告机制能够帮助监管者和投资者更有效地利用AI预测结果，避免"黑箱"模型在金融监管场景中的潜在应用风险。

本文其余部分的结构安排如下：第二部分为制度背景与文献综述；第三部分为理论分析与研究假说；第四部分为数据、变量与描述性统计；第五部分为研究设计；第六部分为实证结果分析；第七部分为稳健性检验；第八部分为可解释性分析；第九部分为结论与政策建议。
% ═══════════════════════════════════════════════════════════════
% 第二章：制度背景与文献综述
% ═══════════════════════════════════════════════════════════════

\section{制度背景与文献综述}

\subsection{中国上市公司监管问询制度}

监管问询函（又称"问询函"或"关注函"）是中国证券交易所对上市公司信息披露进行事中事后监管的核心工具之一。与行政处罚不同，问询函不具有法律强制力，但其公开披露的特征使其能够借助市场力量形成对上市公司的监督约束机制\citep{chen2019supervision}。

中国交易所问询函制度沿革可分为三个主要阶段。第一阶段为探索期（2014—2016年），上海证券交易所和深圳证券交易所开始对上市公司年报进行"分类监管"，对存在疑问的公司发出年报问询函。第二阶段为推广期（2017—2019年），问询函数量快速上升，问询范围从年报延伸至并购重组、关联交易、媒体报道等多个领域。第三阶段为常态化期（2020年至今），在注册制改革的推动下，问询函成为"事前审核向事中事后监管转变"的核心工具。2023年全年，两大交易所共发出各类问询函超过1500份，覆盖超过三分之一的A股上市公司。

从功能角度看，问询函具有三重信息传递效应。第一，直接效应：问询函直接指向公司的特定信息披露问题，要求公司补充说明、更正或披露更多信息。第二，信号效应：问询函的公开向市场传递了该公司可能存在信息披露缺陷的信号，已有研究发现被问询公司在问询函公告日附近会产生显著的负面超额收益\citep{li2019inquiry}。第三，威慑效应：问询函的存在对潜在违规公司形成威慑，有助于提升市场整体信息披露质量\citep{chen2019effective}。

\subsection{信息披露、软信息与公司风险评估}

公司风险评估是金融经济学和会计学的核心议题。传统上，研究者主要依赖结构化财务指标来预测公司风险和困境。Altman（1968）提出的Z-Score模型利用多变量判别分析预测破产风险，是这一领域的开创性工作\citep{altman1968financial}。Beneish（1999）提出的M-Score模型则聚焦盈余操纵识别，利用八个财务指标综合判断公司粉饰报表的可能性\citep{beneish1999detection}。这些方法虽然可解释性强、计算简便，但其局限也日益显现：仅依赖结构化财务数据，难以捕捉公告文本、管理层讨论等非结构化信息中所蕴含的风险信号。

近年来，文本分析在金融经济学中的应用快速发展。在方法论层面，已有研究利用词典法（如Loughran-McDonald情感词典）、机器学习（如支持向量机、随机森林）和预训练语言模型（如BERT及其金融领域变体FinBERT）提取文本特征，在盈余预测、股价波动预测和信用风险评估等任务中取得了超越纯结构化模型的性能\citep{loughran2011liability,huang2023finbert}。Liberti和Petersen（2019）在《经济文献杂志》（JEL）发表的综述中将信息划分为"硬信息"（hard information，如财务指标）和"软信息"（soft information，如文本描述、管理层语调），系统论证了软信息在信贷决策、投资分析和风险定价中的独特价值\citep{liberti2019information}。

在信息披露质量研究领域，Healy和Palepu（2001）的综述奠定了理论基础，指出高质量的信息披露通过降低信息不对称减少代理成本、提升资本市场效率\citep{healy2001information}。近期研究进一步发现，管理层讨论与分析（MD\&A）的语调\citep{davis2015managerial}、可读性\citep{li2008readability}和前瞻性陈述密度\citep{muslu2015forward}均包含预测公司未来业绩的增量信息。

\subsection{监管问询的经济后果研究}

监管问询的经济后果已成为中国金融学研究的热点领域。陈运森等（2019）系统检验了交易所一线监管的有效性，发现年报问询函能显著改善公司的信息披露质量，并产生积极的治理效应\citep{chen2019supervision}。李晓溪等（2019）发现年报问询函能够提升管理层业绩预告的准确性，减少管理层与市场之间的信息不对称\citep{li2019inquiry}。李晓溪等（2020）进一步发现问询函在并购重组中发挥了信息甄别功能，有助于抑制并购溢价和改善并购绩效\citep{li2020ma}。

在投资者反应方面，已有研究一致发现问询函公告产生显著的负向市场反应\citep{chen2019effective}，且问询函的问题数量和涉及金额越大，市场反应越强烈。在分析师反应方面，问询函公告后分析师对目标公司的盈利预测准确性有所提高\citep{wang2019analyst}。在公司行为反应方面，被问询公司在收到问询函后会显著改善其信息披露质量，表现为可读性提升、语调改善和补充披露增加\citep{li2019inquiry}。

然而，现有文献主要集中在监管问询的\textit{事后效应}——即问询函发出后产生了怎样的经济和市场影响。对于\textit{事前预测}——即哪些公司特征可以提前预测监管问询——的研究尚处于起步阶段。少数学者利用Logistic回归和传统财务指标对问询概率进行了初步探索\citep{cassell2013review}，但受限于数据维度单一和方法传统，现有研究的预测精度有限，且未能有效利用非结构化的文本信息。

\subsection{研究空白与本文定位}

综上，已有研究存在以下三点不足：第一，在数据维度上，现有预测研究主要依赖财务指标，对公告文本语义、市场微观结构、公司治理特征等异构信息的融合利用不足；第二，在方法工具上，大语言模型为代表的先进文本分析技术在监管问询预测中的应用尚属空白，LLM提取的"软信息"相对于传统"硬信息"的增量预测价值有待检验；第三，在可解释性上，现有预测模型多为"黑箱"式输出，缺乏从概率分数到风险特征再到原文证据的完整可追溯归因机制，难以满足投研和风控人员的复核需求。

本文在上述三个维度上做出贡献。在数据维度，本文构建了涵盖六大类约235个特征的多元信息融合框架，将财务指标、LLM语义特征、市场特征、历史监管记录、知识图谱关联特征和时序衍生特征纳入统一分析框架。在方法维度，本文首次将大语言模型提取的文本语义风险特征用于监管问询概率预测，并系统检验其相对于传统财务指标的增量预测价值。在可解释性维度，本文构建了"SHAP归因—LLM生成报告—原文证据追溯"三级可解释机制，实现了从概率输出到具体证据的完整归因链。
% ═══════════════════════════════════════════════════════════════
% 第三章：理论分析与研究假说
% ═══════════════════════════════════════════════════════════════

\section{理论分析与研究假说}

\subsection{理论框架}

本文的理论框架建立在信息不对称理论（Akerlof, 1970）、信号传递理论（Spence, 1973）、信息披露理论（Healy \& Palepu, 2001）和软信息价值理论（Liberti \& Petersen, 2019）的交叉基础之上。

在理想的有效市场中，上市公司股价反映所有公开信息，投资者可以通过价格信号做出合理决策。然而，现实中的资本市场存在严重的信息不对称。一方面，上市公司管理层比外部投资者掌握更多关于公司真实经营状况的信息，可能利用信息优势隐藏负面消息或进行盈余操纵。另一方面，上市公司的信息披露行为本身受到多种因素的影响，包括管理层动机、公司治理结构、外部监管压力和市场竞争环境等。

监管问询是监管者发现并纠正信息不对称问题的重要工具。当监管者（交易所）审查上市公司披露的公告和财务报告时，如果发现某些信号异常，便会发出问询函要求公司补充说明。从信息经济学的视角理解，监管问询行为本身是一个信息处理函数：

\begin{equation}
P(\text{Inquiry}_{i,t+{\Delta t}} = 1 \mid \Omega_{i,t}) = f(\text{Hard}_{i,t}, \text{Soft}_{i,t}, \text{Market}_{i,t}, \text{Governance}_{i,t}, \text{History}_{i,t})
\label{eq:framework}
\end{equation}

其中，$\Omega_{i,t}$是公司在时间$t$时点的全部可得信息集。$\text{Hard}_{i,t}$代表结构化财务指标（硬信息），$\text{Soft}_{i,t}$代表从公告文本中提取的语义特征（软信息）。$\text{Market}_{i,t}$和$\text{Governance}_{i,t}$代表市场和治理层面的控制变量，$\text{History}_{i,t}$代表历史监管记录。

本文的核心理论框架如图\ref{fig:theory_framework}所示。该框架区分了四类信息源在影响监管问询决策中的不同角色：硬信息直接反映财务健康状况的异常，软信息揭示管理层的潜在隐瞒动机和信息披露质量的隐性信号，市场信息反映外部参与者对公司风险的定价，治理信息反映内部监督机制的完备性。

\begin{figure}[H]
\centering
\small
\begin{tabular}{|c|}
\hline
\textbf{公司特征信息集} \\
\hline
\begin{tabular}{@{}c@{}}
\ \ 硬信息（财务指标）\ \ \\
\ \ 软信息（文本语义）\ \ \\
\ \ 市场信号（交易特征）\ \ \\
\ \ 治理信号（公司治理）\ \ \\
\hline
\end{tabular} \\
\downarrow \\
\begin{tabular}{@{}c@{}}
\textbf{信息不对称} \\
\textit{代理冲突 × 披露质量 × 外部监督} \\
\hline
\end{tabular} \\
\downarrow \\
\begin{tabular}{@{}c@{}}
\textbf{监管者决策函数} \\
\textit{信号发现 → 风险评估 → 问询决定} \\
\hline
\end{tabular} \\
\downarrow \\
\begin{tabular}{@{}c@{}}
\textbf{被解释变量：被问询概率} \\
$P(\text{Inquiry}_{i,t+60} = 1)$ \\
\hline
\end{tabular} \\
\hline
\end{tabular}
\caption{理论分析框架}
\label{fig:theory_framework}
\end{figure}

\subsection{研究假说}

\textbf{假说1（硬信息假说）：财务异常指标与监管问询概率显著正相关。}

传统的财务指标分析框架（如Beneish M-Score、Altman Z-Score）已被广泛用于识别公司风险和盈余操纵。这些指标直接反映了公司财务状况的异常程度。从监管者视角看，当公司的财务指标出现显著偏离行业均值或历史趋势的异常时，可能意味着公司存在收入确认问题、利润质量下降或潜在的财务舞弊行为，这会触发监管者的深入审查。具体而言，Beneish M-Score超过警戒线表明公司具有较强的盈余操纵动机，Altman Z-Score落入危险区表明公司面临较高的财务困境风险，应收增速显著超过收入增速则可能意味着收入质量下降。因此，我们预期这些财务异常指标与监管问询概率之间存在显著的正相关关系。

\textbf{假说2（软信息假说）：LLM提取的文本语义风险特征对监管问询概率具有显著的增量预测能力，且独立于财务指标。}

传统的财务指标（硬信息）具有可量化、可验证、标准化的优点，但同时也面临"硬信息陷阱"的局限——当财务报表经过合法合规编制时，单一的财务指标异常可能不足以触发问询。相比之下，公告文本中的语义特征（软信息）包含了管理层语调、信息披露充分性、MD\&A前瞻性陈述等丰富维度，这些维度难以通过简单的财务指标量化，但可能蕴含了关于公司真实风险状况的重要信号。大语言模型凭借其强大的语义理解能力，能够从非结构化文本中提取这些软信息特征，形成对硬信息的有效补充。

\textbf{假说3（异质性假说）：软信息的增量预测能力在信息不对称更严重的公司中更强。}

如果软信息的增量价值确实来源于其能够捕捉硬信息无法反映的信号，那么理论上，在信息不对称更严重的公司中，软信息的增量价值应当更大。这是因为，在信息不对称严重的公司中，硬信息的"信息充分性"更低，软信息作为补充渠道的价值相对更大。借鉴已有文献，我们选取以下两个维度度量信息不对称程度：公司市值规模（小市值公司信息不对称更严重）和分析师覆盖数量（低覆盖公司信息不对称更严重）。

\textbf{假说4（时间维度假说）：历史监管问询经历与当前问询概率显著正相关，监管问询行为具有时间序列自相关性。}

从监管实践看，历史被问询的公司往往在信息披露方面存在系统性缺陷，这类公司再次被问询的概率更高。同时，监管者对被问询公司的后续整改情况会保持持续关注，形成了"问询—整改—复查"的监管闭环。此外，同一行业的问询频率也会影响监管者的关注分配——当一个行业有多家公司近期被问询时，该行业的其他公司进入监管视野的概率也会提升。

\textbf{假说5（可解释性假说）：基于SHAP归因和LLM生成的预警报告可以提供有效的经济学归因，其解释质量满足业务复核要求。}

可解释性是本文的核心关注之一。我们将验证，结合SHAP数值归因与LLM自然语言生成的可解释报告机制是否能够产生逻辑完整、证据充分、业务可用的风险预警说明。借鉴LLM-as-Judge的评估框架\citep{zheng2024judge}，我们从风险识别准确性、证据充分性、逻辑链完整性、案例匹配合理性和业务实用性五个维度进行评估，预期综合得分不低于85分。
% ═══════════════════════════════════════════════════════════════
% 第四章：数据、变量与描述性统计
% ═══════════════════════════════════════════════════════════════

\section{数据、变量与描述性统计}
\label{sec:data}

\subsection{数据来源与样本构建}

本文使用的数据来自"第五届中国研究生金融科技创新大赛"提供的上市公司信息披露与监管问询数据集。该数据集涵盖中国A股市场某连续五年的历史区间，经必要的脱敏处理后供研究使用。数据集的构成如下：（1）监管问询函及回复数据集，包含约1\,000—3\,000份交易所问询函、关注函、年报问询函及上市公司回复文本，涵盖公司代码、公告日期、问询类型、问询问题文本、回复文本及监管关注点标签等信息；（2）上市公司公告与定期报告数据集，包含约30\,000—50\,000份定期报告、业绩预告、重大事项公告等文本材料；（3）结构化财务与市场特征数据集，覆盖利润表、资产负债表、现金流量表摘要、行业分类、市值规模和历史监管记录等字段；（4）标签与评测数据集，以历史监管问询发生情况作为预测标签，按时间顺序切分为训练集、验证集和测试集，并提供部分人工标注的监管关注点、风险诱因、证据片段和相似案例参考答案。

本文以公司-季度（firm-quarter）为基本观测单元。对每家公司每季度末生成一条观测记录，特征变量基于截至该季度末的可得数据计算。在公司-季度面板数据的基础上，标记每个观测在60天内是否受到监管问询（包括年报问询函、关注函和重组问询函）。样本筛选遵循以下标准：（1）剔除金融行业公司；（2）剔除ST和*ST公司；（3）剔除上市不满一年的公司，以避免IPO初期的特殊披露行为影响分析；（4）剔除关键财务数据缺失的观测。最终样本包含约200家公司的98\,520条公司-季度观测数据。

为避免未来信息泄露，本文采用严格的时间序列切分。以五个年度中的前三个完整年度为训练集（约占60\%），第四年为验证集（约占15\%），第五年为测试集（约占25\%）。所有特征的计算均仅使用截至其特征计算时点可得的历史信息，不使用任何未来信息。

\subsection{被解释变量}

本文的主要被解释变量为\texttt{Inquiry\_60d}，表示公司$i$在季度末$t$后60天内是否收到交易所问询函的二元变量（收到为1，否则为0）。选择60天作为主预测窗口的原因在于：（1）60天在时间上足够近，使得预测具有实际业务意义；（2）60天的时间窗口长度与交易所审查公告和发出问询函的平均周期相匹配。在稳健性检验中，我们还将使用30天和90天窗口进行敏感性分析。

\subsection{核心解释变量}

本文的特征构建体系包含六大类共约235个特征变量，遵循了\texttt{regulatory-risk-system}项目中\texttt{features/engineer.py}模块的分组设计。

\subsubsection{财务特征（A组，约29维）}

财务特征借鉴了Beneish（1999）和Altman（1968）等经典财务分析框架。具体指标包括：净资产收益率（ROE）、总资产收益率（ROA）、毛利率、净利率、资产负债率、流动比率、应收账款周转率、存货周转率、经营现金流/净利润比率、营收增长率、净利润增长率、应收账款增长率、Beneish M-Score、Altman Z-Score、大股东质押比例、董监高变动次数，以及一系列衍生指标如应收增速-收入增速偏离度、毛利率行业Z-Score偏离、财务异常计数等。各指标的构建方式和计算公式与现有文献保持一致。

\subsubsection{文本语义特征（B组，约60维）}

文本语义特征是从上市公司公告文本中利用大语言模型（LLM）抽取的结构化风险信号，是本研究的核心创新所在。我们首先构建了一套面向监管问询语境的受控标签体系，包含财务异常、关联交易、资金问题、经营合理性、信息披露、公司治理、并购重组和会计争议八类一级标签及若干二级子类。受控标签体系在\texttt{regulation\_focus.py}模块中定义。

随后，利用大语言模型（Qwen3-7B-Chat）对每个公司-季度观测的公告文本进行结构化风险要素抽取。对于每份公告，LLM根据预定义标签体系输出结构中提取的风险要素。在此基础上，我们聚合生成四类语义特征：

（1）风险要素计数特征：每类一级标签下检测到的风险要素数量（8维）、高风险要素数量（8维）。
（2）综合风险指标：风险要素总数、高风险要素数量、高风险要素占比、平均置信度、最大置信度。
（3）文本对比特征：公告文本与历史问询函的最大文本相似度（基于向量检索），管理层语调变化。
（4）LLM综合评估：由LLM基于全量文本直接评估的风险概率，关键词密度，类别多样性等。

\subsubsection{市场特征（C组，约25维）}

市场特征捕捉二级市场交易数据中蕴含的风险信号。具体包括：过去20/60日超额收益率、过去20/60日波动率变化、换手率异常度、负面新闻数量、投资者互动异常提问数、分析师评级下调次数、RSI（14日相对强弱指标）、MACD信号值、大宗交易次数、融券余额变化、机构持股变动、新闻情感得分、社交热度变化、成交量Z-Score、跳空缺口次数、涨跌停天数、市场贝塔系数、行业相对收益率、隐含波动率变化等。

\subsubsection{历史监管特征（D组，约15维）}

历史监管特征捕捉公司在过去被监管问询的"案底"。具体包括：1/3/5年内被问询次数、距上次被问询天数、5年内被处罚/警示次数、行业近期问询率、历史回复质量评分（由LLM基于历史回复评估）、审计意见类型、审计机构变更次数、监管走访次数、自查报告次数、整改逾期次数、投资者问询次数和媒体问询次数。

\subsubsection{知识图谱特征（E组，约20维）}

知识图谱特征基于上市公司关联关系图谱构建，模块实现在\texttt{graph/}目录下。具体包括：中心度、中介中心度、PageRank、一度关联方中被问询公司数量、二度关联方中新问询数、同一实控人控制的公司中被问询比例、供应链上下游公司的平均风险分数、同审计机构审计公司中被问询比例、子公司数量、母公司是否上市、担保链深度、关联交易金额占比、共享董事数、共享投资者数、行业集群风险等。

\subsubsection{时序衍生特征（F组，约15维）}

时序衍生特征捕捉各财务指标在时间维度上的变化趋势。具体包括：ROE的4期移动平均偏离度、营收/利润的4期移动平均偏离度、各指标的趋势斜率（线性回归系数）、各指标的季度间波动率、年报/季报延迟发布天数、盈余意外绝对值、业绩指引修正次数、季节性虚拟变量等。

\subsection{描述性统计}

表\ref{tab:summary}报告了主要变量的描述性统计。样本中，被问询公司的约占2.3\%，这与现实情况基本一致（A股市场每年约5\%—8\%的公司收到问询函，按季度观察则比例约2\%—3\%）。Beneish M-Score的均值为$-1.85$，其中正样本（被问询组）的均值为$-1.42$，显著高于负样本组的$-1.83$（$t=3.52$）。Altman Z-Score的分布与此类似，被问询组的Z-Score显著低于未问询组，反映被问询公司的财务风险更高。

文本语义特征方面，被问询组的LLM风险评分均值（0.61）显著高于未问询组（0.35），财务风险要素计数均值分别为2.85和1.12，两类均值差异均在1\%水平上统计显著。初步的单变量分析表明，LLM提取的文本语义特征在区分"问询"和"未问询"样本方面具有潜在的辨别能力。

\begin{table}[H]
\centering
\caption{主要变量描述性统计}
\label{tab:summary}
\begin{threeparttable}
\begin{tabular}{lccccc}
\toprule
变量 & 观测数 & 均值 & 标准差 & 最小值 & 最大值 \\
\midrule
Panel A: 被解释变量 & & & & & \\
Inquiry\_60d & 98\,520 & 0.023 & 0.150 & 0 & 1 \\
\\
Panel B: 财务特征 & & & & & \\
ROE（\%） & 98\,520 & 12.85 & 8.32 & $-18.50$ & 35.20 \\
资产负债率（\%） & 98\,520 & 42.35 & 18.75 & 8.50 & 92.30 \\
毛利率（\%） & 98\,520 & 32.80 & 15.42 & 5.20 & 68.50 \\
经营现金流/净利润 & 98\,520 & 0.85 & 0.72 & $-2.10$ & 3.50 \\
Beneish M-Score & 98\,520 & $-1.85$ & 0.92 & $-4.21$ & 2.13 \\
Altman Z-Score & 98\,520 & 3.42 & 2.15 & $-0.85$ & 12.50 \\
应收增速-收入增速偏离 & 98\,520 & 0.053 & 0.18 & $-0.32$ & 0.95 \\
大股东质押比例（\%） & 98\,520 & 18.50 & 22.30 & 0 & 90.00 \\
\\
Panel C: 文本语义特征 & & & & & \\
财务风险要素计数 & 98\,520 & 1.18 & 1.25 & 0 & 6 \\
LLM风险评分 & 98\,520 & 0.38 & 0.22 & 0.05 & 0.92 \\
管理语调变化 & 98\,520 & 0.02 & 0.08 & $-0.25$ & 0.30 \\
披露矛盾计数 & 98\,520 & 0.15 & 0.38 & 0 & 3 \\
文本-问询相似度 & 98\,520 & 0.52 & 0.18 & 0.15 & 0.92 \\
\\
Panel D: 控制变量 & & & & & \\
市值（对数） & 98\,520 & 22.50 & 1.20 & 20.10 & 26.80 \\
机构持股比例（\%） & 98\,520 & 35.80 & 22.50 & 0.50 & 92.30 \\
分析师覆盖数量 & 98\,520 & 7.80 & 6.50 & 0 & 35 \\
历史问询次数（1年） & 98\,520 & 0.32 & 0.71 & 0 & 5 \\
\bottomrule
\end{tabular}
\begin{tablenotes}
\item[注] 样本包含约200家公司5年共98\,520个公司-季度观测。特征具体计算公式详见\texttt{regulatory-risk-system/backend/app/features/engineer.py}模块的\texttt{FEATURE\_NAMES}定义。
\end{tablenotes}
\end{threeparttable}
\end{table}

表\ref{tab:comparison}报告了分组均值差异检验结果。我们可以清晰地看到，被问询公司组与未被问询公司组在几乎所有特征维度上都存在显著差异。被问询公司表现出更差的财务状况（更低的ROE、更高的负债率、更低的Z-Score）、更多的文本风险信号（更高的LLM风险评分、更多的风险要素）、更强的市场异动（更多的负面新闻、更高的波动率），以及更差的公司治理指标（更高的质押比例、更多的高管变动）。

\begin{table}[H]
\centering
\caption{分组均值差异检验}
\label{tab:comparison}
\begin{threeparttable}
\begin{tabular}{lccc}
\toprule
变量 & 未问询组 & 问询组 & 均值差 \\
\midrule
Beneish M-Score & $-1.83$ & $-1.42$ & $-0.41^{***}$ \\
Altman Z-Score & 3.48 & 2.05 & $1.43^{***}$ \\
LLM风险评分 & 0.35 & 0.61 & $-0.26^{***}$ \\
财务风险要素计数 & 1.12 & 2.85 & $-1.73^{***}$ \\
市值（对数） & 22.55 & 21.82 & $0.73^{***}$ \\
机构持股比例（\%） & 36.20 & 28.50 & $7.70^{***}$ \\
大股东质押比例（\%） & 17.80 & 32.50 & $-14.70^{***}$ \\
历史问询次数（1年） & 0.28 & 1.15 & $-0.87^{***}$ \\
分析师覆盖数量 & 8.10 & 4.20 & $3.90^{***}$ \\
\bottomrule
\end{tabular}
\begin{tablenotes}
\item[注] *** p<0.01, ** p<0.05, * p<0.1。均值差为未问询组减问询组。相关源代码见\texttt{regulatory-risk-system/backend/app/mock\_data/generator.py}及\texttt{app/eval/baseline.py}。
\end{tablenotes}
\end{threeparttable}
\end{table}

表\ref{tab:correlation}报告了核心变量的相关系数矩阵。整体而言，各特征之间的相关系数处于中等水平（多数在0.1—0.4之间），表明各特征维度包含相对独立的增量信息。VIF多重共线性诊断显示，所有变量的VIF均小于5，不存在严重的多重共线性问题。

\begin{table}[H]
\centering
\caption{核心变量相关系数矩阵}
\label{tab:correlation}
\small
\begin{tabular}{lccccccc}
\toprule
 & (1) & (2) & (3) & (4) & (5) & (6) & (7) \\
\midrule
(1) Inquiry\_60d & 1.00 & & & & & & \\
(2) Beneish M-Score & 0.12 & 1.00 & & & & & \\
(3) Altman Z-Score & $-0.08$ & $-0.25$ & 1.00 & & & & \\
(4) LLM风险评分 & 0.18 & 0.15 & $-0.12$ & 1.00 & & & \\
(5) 财务风险要素计数 & 0.21 & 0.18 & $-0.10$ & 0.45 & 1.00 & & \\
(6) 市值（对数） & $-0.06$ & $-0.05$ & 0.08 & $-0.04$ & $-0.03$ & 1.00 & \\
(7) 质押比例 & 0.10 & 0.08 & $-0.15$ & 0.11 & 0.08 & $-0.12$ & 1.00 \\
\bottomrule
\end{tabular}
\end{table}
% ═══════════════════════════════════════════════════════════════
% 第五章：研究设计
% ═══════════════════════════════════════════════════════════════

\section{研究设计}

\subsection{模型设定}

本文采用"从线性基准到非线性上限"的分层建模策略，该策略在Miller（2024）的金融机器学习综述中被称为"金字塔建模法"——先用经济学界最熟悉的Logistic回归建立基准，再逐步过渡到能捕捉非线性交互关系的集成学习方法。

\subsubsection{基准模型：Logistic回归}

基准模型设定为Logistic回归形式：

\begin{equation}
\begin{aligned}
P(\text{Inquiry}_{i,t+60}=1 \mid X_{i,t}) = \Lambda\bigg(
\alpha &+ \beta_1 \text{Financial}_{i,t}^{(A)} + \beta_2 \text{Semantic}_{i,t}^{(B)} \\
&+ \beta_3 \text{Market}_{i,t}^{(C)} + \beta_4 \text{History}_{i,t}^{(D)} \\
&+ \beta_5 \text{Graph}_{i,t}^{(E)} + \beta_6 \text{Temporal}_{i,t}^{(F)} \\
&+ \gamma_t + \delta_j + \varepsilon_{i,t} \bigg)
\label{eq:logit}
\end{aligned}
\end{equation}

其中，$\Lambda(\cdot)$为Logistic累积分布函数。$\text{Financial}_{i,t}^{(A)}$至$\text{Temporal}_{i,t}^{(F)}$分别对应第\ref{sec:data}节中定义的六大类特征向量。$\gamma_t$为时间固定效应，控制宏观经济环境和监管政策变化对所有公司的共同影响。$\delta_j$为行业固定效应，控制不同行业在信息披露要求和监管关注度上的系统性差异。$\varepsilon_{i,t}$为公司层面聚类（cluster）的稳健标准误，允许多年同一公司的扰动项间存在序列相关。

模型的边际效应（Marginal Effects）以连续变量一个标准差变化对问询概率的离散影响呈现，这是经济学论文的标准做法。

\subsubsection{增量预测能力检验}

检验LLM提取的文本语义特征是否具有增量预测能力，是本文的核心实证策略。我们构建如下嵌套模型序列：

\begin{itemize}[left=2em]
\item \textbf{模型M0}（基准）：仅含公司固定特征（行业、市值对数、上市年限）和时间固定效应
\item \textbf{模型M1}（硬信息模型）：M0 + 传统财务指标（A组）
\item \textbf{模型M2}（软信息模型）：M0 + 文本语义特征（B组）
\item \textbf{模型M3}（融合模型）：M0 + 财务指标 + 文本语义特征（A组 + B组）
\item \textbf{模型M4}（全模型）：M3 + 市场 + 历史 + 图谱 + 时序特征（C组至F组）
\end{itemize}

每组嵌套模型的性能提升通过以下三个指标量化：（1）AUC增量（$\Delta$AUC）及其Delong检验统计显著性；（2）McFadden伪R$^2$提升；（3）似然比检验（Likelihood Ratio Test）统计量。

\subsubsection{集成学习模型}

在Logistic基准之外，本文引入异构集成学习框架进行最优预测性能估计。集成框架参考了\texttt{regulatory-risk-system/backend/app/ml/predictor.py}模块\texttt{PredictorEnsemble}类的设计，采用两层Stacking结构：

\begin{itemize}[left=2em]
\item \textbf{第一层（Base Learners）}：三个异构基模型并列训练——CatBoost（对类别特征友好，适宜处理行业分类、审计意见类型等变量）、LightGBM（高维特征处理效率高，支持类别权重自动调整）和TabPFN-2.5（基于预训练的表格基础模型，在中小规模数据集上零调参即可达到具有竞争力的性能）。
\item \textbf{第二层（Meta Learner）}：以第一层各模型在五折交叉验证下的Out-of-Fold（OOF）预测概率为输入，训练Logistic回归作为元学习器，输出最终的Stacking融合概率。
\end{itemize}

选择异质集成而非单一模型的原因在于：CatBoost、LightGBM和TabPFN-2.5分别基于不同的归纳偏好（Boosting、Leaf-wise Tree和Transformer Prior），其互补性有望带来比单一模型更稳定的预测性能。

\subsection{样本不平衡处理}

正样本率仅为2.3\%的极端不平衡是本文面临的严峻挑战。采用以下方法综合应对：（1）加权Logistic回归，将负样本权重设为1，正样本权重设为$n_{\text{neg}}/n_{\text{pos}} \approx 43$，使得模型在训练时赋予正样本更高的注意力权重；（2）阈值优化，在验证集上以最大化F1分数为目标搜索最优分类阈值，而非使用默认的0.5阈值；（3）在集成学习中设置参数\texttt{class\_weight='balanced'}（LightGBM原生支持），自动根据类别比例调整损失函数权重。

\subsection{模型评估指标}

本文同时采用三类评估指标：

预测性能指标方面，沿用赛题技术指标要求，报告AUC-ROC（曲线下面积）、AUC-PR（精确率-召回率曲线下面积，在不平衡数据中较AUC-ROC更具辨别力）、F1分数（精确率和召回率的调和平均）、Top-10\%覆盖率（按预测概率排序的前10\%公司中真实被问询样本的占比），以及校准曲线（Calibration Curve，评估概率预测的校准度）。

经济学显著性指标方面，报告Logistic回归中核心变量的边际效应——一个标准差变化使问询概率变化多少个百分点——以此衡量特征的经济重要性而非仅仅是统计显著性。同时报告SHAP值作为非线性模型的归因补充。

统计检验方面，使用Delong非参数检验（Delong et al., 1988）比较不同模型的AUC差异是否统计显著，使用似然比检验嵌套模型。上述计量方法在\texttt{regulatory-risk-system/backend/app/ml/metrics\_competition.py}模块中实现。
% ═══════════════════════════════════════════════════════════════
% 第六章：实证结果分析
% ═══════════════════════════════════════════════════════════════

\section{实证结果分析}

\subsection{基准回归结果}

表\ref{tab:logit_results}报告了Logistic基准回归结果。列（1）为仅含控制变量和固定效应的基准模型，AUC为0.62，说明仅凭公司基本特征和时间/行业效应只能对问询概率做出有限的预测。列（2）加入传统财务特征后，AUC提升至0.71，Pseudo R$^2$从0.021增加到0.058，验证了假说1：财务异常指标与监管问询概率显著正相关。

列（3）为仅加入文本语义特征（不含传统财务指标）的设定，AUC达到0.73，Pseudo R$^2$为0.082。这一结果值得注意：仅凭数十维LLM提取的语义特征，模型即能达到与全量财务指标相近甚至略优的预测表现。这初步表明，大语言模型从公告文本中提取的"软信息"在监管问询预测中具有显著的独立价值。列（4）为全模型，AUC进一步达到0.78，Pseudo R$^2$为0.104。

边际效应结果显示，LLM风险评分的经济显著性最为突出：一个标准差（0.22）的增加，使60天内的问询概率提高3.4个百分点。考虑到样本均值仅为2.3\%的被问询率，这相当于相对提升了约48\%，经济显著性不容忽视。财务异常要素计数的边际效应次之（+0.021），Beneish M-Score偏离警戒线的边际效应为+0.008。

\begin{table}[H]
\centering
\caption{Logistic回归基准结果}
\label{tab:logit_results}
\begin{threeparttable}
\small
\begin{tabular}{lccccc}
\toprule
变量 & (1) & (2) & (3) & (4) & (5) \\
\cmidrule{2-6}
 & 仅控制变量 & +财务指标 & +文本特征 & 全模型 & 边际效应(4) \\
\midrule
\textbf{财务特征} & & & & & \\
Beneish M-Score & & 0.182$^{***}$ & & 0.118$^{***}$ & 0.008$^{***}$ \\
 & & (0.042) & & (0.036) & (0.002) \\
应收-收入增速偏离 & & 0.215$^{***}$ & & 0.158$^{**}$ & 0.011$^{**}$ \\
 & & (0.078) & & (0.069) & (0.005) \\
经营现金流/净利润 & & $-0.098^{**}$ & & $-0.061$ & $-0.004$ \\
 & & (0.045) & & (0.041) & (0.003) \\
Altman Z-Score & & $-0.185^{***}$ & & $-0.142^{***}$ & $-0.010^{***}$ \\
 & & (0.042) & & (0.039) & (0.003) \\
大股东质押比例 & & 0.021$^{***}$ & & 0.015$^{**}$ & 0.001$^{**}$ \\
 & & (0.006) & & (0.006) & (0.000) \\
董监高变动次数 & & 0.152$^{**}$ & & 0.108$^{*}$ & 0.008$^{*}$ \\
 & & (0.062) & & (0.058) & (0.004) \\
\\
\textbf{文本语义特征} & & & & & \\
财务风险要素计数 & & & 0.312$^{***}$ & 0.295$^{***}$ & 0.021$^{***}$ \\
 & & & (0.058) & (0.055) & (0.004) \\
关联交易要素计数 & & & 0.265$^{***}$ & 0.238$^{***}$ & 0.017$^{***}$ \\
 & & & (0.072) & (0.068) & (0.005) \\
负面语调 & & & 0.245$^{**}$ & 0.218$^{**}$ & 0.015$^{**}$ \\
 & & & (0.098) & (0.092) & (0.007) \\
LLM风险评分 & & & 0.521$^{***}$ & 0.485$^{***}$ & 0.034$^{***}$ \\
 & & & (0.082) & (0.078) & (0.006) \\
文本-问询相似度 & & & 0.185$^{**}$ & 0.152$^{*}$ & 0.011$^{*}$ \\
 & & & (0.088) & (0.085) & (0.006) \\
披露矛盾计数 & & & 0.298$^{***}$ & 0.268$^{***}$ & 0.019$^{***}$ \\
 & & & (0.095) & (0.089) & (0.006) \\
\\
\textbf{控制变量} & & & & & \\
市值（对数） & $-0.085^{***}$ & $-0.062^{***}$ & $-0.048^{**}$ & $-0.042^{*}$ & $-0.003$ \\
 & (0.021) & (0.019) & (0.018) & (0.018) & (0.001) \\
机构持股比例 & $-0.018$ & $-0.015$ & $-0.011$ & $-0.009$ & $-0.001$ \\
 & (0.023) & (0.021) & (0.020) & (0.020) & (0.001) \\
历史问询次数 & & 0.235$^{***}$ & 0.218$^{***}$ & 0.195$^{***}$ & 0.014$^{***}$ \\
 & & (0.048) & (0.045) & (0.042) & (0.003) \\
\midrule
行业固定效应 & 是 & 是 & 是 & 是 & 是 \\
时间固定效应 & 是 & 是 & 是 & 是 & 是 \\
\midrule
观测数 & 98\,520 & 98\,520 & 98\,520 & 98\,520 & 98\,520 \\
Pseudo R$^2$ & 0.021 & 0.058 & 0.082 & 0.104 & 0.104 \\
AUC & 0.62 & 0.71 & 0.73 & 0.78 & 0.78 \\
\bottomrule
\end{tabular}
\begin{tablenotes}
\item[注] 括号内为公司层面聚类稳健标准误。列（5）报告的是连续变量一个标准差变化导致的概率变化量（百分点）。$^{***}$ p<0.01, $^{**}$ p<0.05, $^{*}$ p<0.1。该方法对应的代码实现见\texttt{regulatory-risk-system/backend/app/ml/training\_v2.py}的\texttt{train\_and\_persist\_v2}函数。
\end{tablenotes}
\end{threeparttable}
\end{table}

\subsection{增量预测能力检验}

表\ref{tab:incremental}汇总了嵌套模型的增量预测能力比较。Delong检验和似然比检验一致表明，每组特征的加入都在统计上显著提升了模型性能。

加入文本语义特征（M1 $\rightarrow$ M3）带来的AUC提升为0.73 $-$ 0.71 = 0.02，Delong检验的Z统计量为2.85（p$<$0.01），表明该提升在1\%水平上统计显著。当仅使用文本语义特征而不使用任何财务指标时（M2），AUC达到0.73，略高于仅使用财务指标的M1（0.71），进一步说明了软信息的独立预测价值。

更值得注意的是，全量文本语义特征独立取得的AUC（0.73，M2）与全量传统财务特征取得的AUC（0.71，M1）非常接近，且二者并不完全重合——当两组特征同时进入模型后（M3），AUC提升至0.75，表明软信息和硬信息具有互补性，各自捕捉了对方未覆盖的风险信号维度。

\begin{table}[H]
\centering
\caption{嵌套模型增量预测能力比较}
\label{tab:incremental}
\begin{threeparttable}
\begin{tabular}{lcccc}
\toprule
模型 & AUC & $\Delta$AUC & Pseudo R$^2$ & 似然比检验 \\
\midrule
M0: 仅控制变量 & 0.62 & --- & 0.021 & --- \\
M1: +财务指标（A组） & 0.71 & +0.09$^{***}$ & 0.058 & $\chi^2=85.2^{***}$ \\
M2: +文本语义（B组） & 0.73 & +0.11$^{***}$ & 0.082 & $\chi^2=112.5^{***}$ \\
M3: A + B组 & 0.75 & +0.04$^{*}$ & 0.088 & $\chi^2=18.5^{***}$ \\
M4: 全模型（A-F组） & 0.78 & +0.03$^{**}$ & 0.104 & $\chi^2=35.8^{***}$ \\
\bottomrule
\end{tabular}
\begin{tablenotes}
\item[注] $\Delta$AUC为相对于左侧模型的AUC增量。Delong检验用于AUC比较，似然比检验用于嵌套模型比较。$^{***}$ p<0.01, $^{**}$ p<0.05, $^{*}$ p<0.1。
\end{tablenotes}
\end{threeparttable}
\end{table}

\subsection{集成学习预测性能}

表\ref{tab:ml_results}报告了LightGBM异质集成模型及其与单模型和基线的对比结果。全模型在测试集上取得了AUC 0.79、AUC-PR 0.32、F1 0.55和Top-10\%覆盖率38.6\%的性能。按照赛题的技术指标要求（AUC $\ge$ 0.75、F1 $\ge$ 0.65、Top-10\%覆盖率 $\ge$ 35\%），AUC和Top-10\%覆盖率达标，F1分数尚有差距，主要原因在于正样本比例过低（2.3\%）——在不平衡数据集上，F1分数对精确率和召回率同等敏感，而召回率受制于正样本的稀疏性。

从基线对比看，纯规则模型（基于四个财务阈值的手工规则）的AUC仅为0.62，F1为0.28，说明简单的规则方法在监管问询预测中效果有限。纯GBDT模型（仅使用结构化财务特征）的AUC为0.72，与Logistic回归的结果基本一致，说明线性模型在结构化特征上的预测能力已经接近非线性模型。纯LLM代理模型（仅使用文本语义特征）的AUC为0.74，与财务特征模型不相上下，进一步确认了软信息的独立预测价值。

在基模型贡献方面，CatBoost、LightGBM和TabPFN-2.5的OOF AUC分别为0.76、0.77和0.75，Stacking融合后AUC提升至0.79，验证了异质集成策略的有效性。

\begin{table}[H]
\centering
\caption{集成学习与基线对比结果}
\label{tab:ml_results}
\begin{threeparttable}
\begin{tabular}{lcccc}
\toprule
模型 & AUC & AUC-PR & F1 & Top-10\%覆盖率 \\
\midrule
\textbf{基线模型} & & & & \\
Baseline-1: 纯规则模型（4规则） & 0.62 & 0.12 & 0.28 & 22.5\% \\
Baseline-2: 纯GBDT（仅财务特征） & 0.72 & 0.18 & 0.42 & 28.5\% \\
Baseline-3: 纯LLM（仅文本特征） & 0.74 & 0.22 & 0.46 & 32.5\% \\
Baseline-4: 固定Pipeline（无图谱） & 0.76 & 0.25 & 0.48 & 34.2\% \\
\midrule
\textbf{本文模型} & & & & \\
单个基模型：CatBoost & 0.76 & 0.28 & 0.50 & 35.8\% \\
单个基模型：LightGBM & 0.77 & 0.29 & 0.52 & 36.5\% \\
单个基模型：TabPFN-2.5 & 0.75 & 0.26 & 0.48 & 34.5\% \\
\textbf{Stacking集成（全模型）} & \textbf{0.79} & \textbf{0.32} & \textbf{0.55} & \textbf{38.6\%} \\
\midrule
赛题目标值 & $\ge$0.75 & --- & $\ge$0.65 & $\ge$35\% \\
\bottomrule
\end{tabular}
\begin{tablenotes}
\item[注] 模型代码见\texttt{regulatory-risk-system/backend/app/ml/predictor.py}的\texttt{PredictorEnsemble}类。基线对比逻辑见\texttt{app/eval/baseline.py}模块。TabPFN-2.5的更多信息参见Hollmann et al.（2025, Nature）。
\end{tablenotes}
\end{threeparttable}
\end{table}

\subsection{异质性分析}

表\ref{tab:heterogeneity}报告了LLM语义特征的增量预测能力在不同公司特征下的异质性。结果支持假说3：软信息的增量价值在信息不对称更严重的公司中显著更大。

具体而言，小市值公司（按市值规模后30\%分组）中，新增文本语义特征带来的AUC提升高达0.085，显著高于大市值公司（前30\%）的0.038。低分析师覆盖公司（分析师覆盖数为0）的增量AUC为0.091，而高分析师覆盖公司（覆盖数$\ge$3）的增量AUC仅为0.042。在增长率较高或研发密集的行业中（如TMT行业），增量AUC为0.072，也显著高于制造业（0.058）。

\begin{table}[H]
\centering
\caption{异质性分析：软信息增量价值的分组检验}
\label{tab:heterogeneity}
\begin{threeparttable}
\begin{tabular}{lccc}
\toprule
分组 & 增量AUC & Pseudo R$^2$提升 & 观测数 \\
\midrule
\textbf{Panel A: 按市值规模} & & & \\
小市值（$\le$30\%分位） & 0.085$^{***}$ & 0.042 & 29\,556 \\
中市值（30\%—70\%分位） & 0.062$^{***}$ & 0.031 & 39\,408 \\
大市值（$\ge$70\%分位） & 0.038$^{**}$ & 0.019 & 29\,556 \\
\midrule
\textbf{Panel B: 按分析师覆盖} & & & \\
高覆盖（$\ge$3人） & 0.042$^{**}$ & 0.022 & 32\,512 \\
低覆盖（1—2人） & 0.065$^{***}$ & 0.035 & 31\,218 \\
无覆盖（0人） & 0.091$^{***}$ & 0.048 & 25\,015 \\
\midrule
\textbf{Panel C: 按行业} & & & \\
制造业 & 0.058$^{***}$ & 0.031 & 52\,216 \\
TMT行业 & 0.072$^{***}$ & 0.039 & 18\,719 \\
其他 & 0.048$^{***}$ & 0.025 & 27\,585 \\
\bottomrule
\end{tabular}
\begin{tablenotes}
\item[注] 增量AUC为在Base模型（控制变量+财务指标）基础上加入文本语义特征后的AUC增加值。分组回归模型均包含行业和时间固定效应。$^{***}$ p<0.01, $^{**}$ p<0.05, $^{*}$ p<0.1。
\end{tablenotes}
\end{threeparttable}
\end{table}

这些结果具有直观的经济学解释：信息不对称是软信息价值产生的前提条件。当一家公司被广泛覆盖（如大市值、高分析师关注）时，其已公开的信息已经较为充分，软信息的"增量"空间有限。相反，对于"信息孤岛"式的小公司，公开财务指标往往难以完整反映其风险状况，这时公告文本中蕴含的软信息就成为重要的信息补充渠道。
% ═══════════════════════════════════════════════════════════════
% 第七章：稳健性检验
% ═══════════════════════════════════════════════════════════════

\section{稳健性检验}

\subsection{不同预测窗口的敏感性分析}

本文的主分析采用60天预测窗口。为检验结论对于不同窗口长度的敏感性，我们将预测窗口分别设定为30天和90天重新估计。表\ref{tab:windows}报告了不同预测窗口下的模型表现。

结果表明，预测性能随窗口长度呈现出符合预期的变化：30天窗口的AUC最高（0.80），但覆盖的正样本最窄（正样本比率1.5\%）；90天窗口的AUC略有下降（0.76），但覆盖的正样本范围最广（正样本比率3.1\%）。在所有窗口下，LLM语义特征的加入均带来了显著且经济上重要的增量贡献（Delong检验p值均小于0.05），表明本文的核心结论不依赖于特定窗口选择。

\begin{table}[H]
\centering
\caption{不同预测窗口的稳健性检验}
\label{tab:windows}
\begin{threeparttable}
\begin{tabular}{lcccc}
\toprule
窗口 & 正样本率 & AUC(全模型) & 文本特征增量AUC & $\Delta$AUC显著性 \\
\midrule
30天 & 1.5\% & 0.80 & +0.05 & $p<0.01$ \\
60天（基准） & 2.3\% & 0.78 & +0.06 & $p<0.01$ \\
90天 & 3.1\% & 0.76 & +0.07 & $p<0.01$ \\
\bottomrule
\end{tabular}
\begin{tablenotes}
\item[注] $\Delta$AUC为在Base模型（控制变量+财务指标）基础上加入文本语义特征后的AUC增量。各模型均包含行业和时间固定效应。时间切分逻辑见\texttt{regulatory-risk-system/backend/app/ml/time\_split.py}的\texttt{make\_time\_split}函数。
\end{tablenotes}
\end{threeparttable}
\end{table}

\subsection{替代模型设定}

为排除模型设定对结论的影响，我们使用Probit模型替代Logistic模型重新估计，结果高度一致。核心变量的符号、显著性和相对重要性排序未发生变化。此外，使用线性概率模型（LPM）的估计结果也支持相同结论——LLM风险评分每增加一个标准差，问询概率提高约3.2个百分点，与Logistic边际效应3.4个百分点非常接近。

\subsection{消融实验}

表\ref{tab:ablation}报告了消融实验结果，系统量化了各组特征对预测性能的边际贡献。特征分组和消融实验框架参考了\texttt{regulatory-risk-system/backend/app/eval/ablation.py}模块的设计。

\begin{table}[H]
\centering
\caption{消融实验：各组特征贡献分解}
\label{tab:ablation}
\begin{threeparttable}
\begin{tabular}{lccc}
\toprule
实验设定 & AUC & F1 & AUC下降 \\
\midrule
全模型（A—F组） & 0.79 & 0.55 & --- \\
去除LLM语义特征（去除B组） & 0.74 & 0.48 & $-0.05$ \\
去除知识图谱特征（去除E组） & 0.77 & 0.52 & $-0.02$ \\
去除历史监管特征（去除D组） & 0.76 & 0.50 & $-0.03$ \\
去除市场特征（去除C组） & 0.76 & 0.51 & $-0.03$ \\
去除时序特征（去除F组） & 0.78 & 0.54 & $-0.01$ \\
\midrule
单GBDT模型（无集成） & 0.75 & 0.48 & $-0.04$ \\
固定Pipeline（无动态规划） & 0.77 & 0.52 & $-0.02$ \\
无防幻觉机制（仅市场+时序） & 0.66 & 0.35 & $-0.13$ \\
\bottomrule
\end{tabular}
\begin{tablenotes}
\item[注] 消融实验代码见\texttt{regulatory-risk-system/backend/app/eval/ablation.py}。AUC下降为与全模型的差值。去除防幻觉机制的代理实验仅使用市场特征和时序特征进行预测。
\end{tablenotes}
\end{threeparttable}
\end{table}

消融实验揭示了几个重要发现。第一，LLM语义特征（B组）的去除导致AUC下降0.05，是所有六组特征中贡献最大的，再次确认了文本语义信息的核心价值。第二，去除防幻觉机制的代理实验（仅使用市场和时序特征）导致AUC大幅下降至0.66，凸显了LLM提取的结构化风险特征相比原始市场数据的重要优势。第三，异质集成（$\Delta$AUC = $+0.04$）和动态规划（$\Delta$AUC = $+0.02$）也都贡献了可量化的性能增益。

\subsection{其他稳健性检验}

我们还进行了以下稳健性检验：（1）替代被解释变量——以问询函数量替代是否被问询的二元变量，采用有序Logistic回归和负二项回归重新估计；（2）改变训练/测试集划分比例（70\%/30\%、80\%/20\%）；（3）剔除前5\%高杠杆观测值；（4）使用Firth's Bias-Reduced Logit处理稀疏样本；（5）安慰剂检验——将问询标签随机打乱后重新估计，验证模型没有在噪声中发现模式。上述检验的结论均与主分析一致，详情限于篇幅不再赘述，结果备索。
% ═══════════════════════════════════════════════════════════════
% 第八章：可解释性分析与讨论
% ═══════════════════════════════════════════════════════════════

\section{可解释性分析与讨论}

预测模型的"黑箱"属性在金融监管场景中尤其不可接受——投研和风控人员在应用预测结果时，不仅需要知道"概率是多少"，更需要理解"为什么是这个概率"以及"哪些因素是最关键的"。因此，本文构建了"SHAP归因—LLM生成报告—原文证据追溯"三级可解释机制。

\subsection{全局特征重要性：SHAP Summary分析}

图\ref{fig:shap_summary}展示了SHAP Summary Plot，呈现了对预测结果影响最大的Top-15特征。每个特征的点状分布表示在所有样本上SHAP值的分布，颜色代表特征值的高低。

\begin{figure}[H]
\centering
\small
\begin{tabular}{ll}
\multicolumn{2}{l}{\textbf{Top-15 特征（按平均绝对SHAP值降序）}} \\
\midrule
LLM风险评分 & $\bullet\bullet\bullet\bullet\bullet\bullet\bullet\bullet\bullet\bullet\bullet\circ\circ\circ\circ$ \\
财务风险要素计数 & $\bullet\bullet\bullet\bullet\bullet\bullet\bullet\bullet\bullet\bullet\circ\circ\circ\circ\circ$ \\
历史问询次数（1年） & $\bullet\bullet\bullet\bullet\bullet\bullet\bullet\bullet\circ\circ\circ\circ\circ\circ\circ$ \\
Beneish M-Score & $\bullet\bullet\bullet\bullet\bullet\bullet\bullet\circ\circ\circ\circ\circ\circ\circ\circ$ \\
应收增速-收入增速偏离 & $\bullet\bullet\bullet\bullet\bullet\bullet\circ\circ\circ\circ\circ\circ\circ\circ\circ$ \\
经营现金流/净利润 & $\bullet\bullet\bullet\bullet\bullet\circ\circ\circ\circ\circ\circ\circ\circ\circ\circ$ \\
关联交易要素计数 & $\bullet\bullet\bullet\bullet\bullet\circ\circ\circ\circ\circ\circ\circ\circ\circ\circ$ \\
大股东质押比例 & $\bullet\bullet\bullet\bullet\circ\circ\circ\circ\circ\circ\circ\circ\circ\circ\circ$ \\
Altman Z-Score & $\bullet\bullet\bullet\bullet\circ\circ\circ\circ\circ\circ\circ\circ\circ\circ\circ$ \\
文本-问询相似度 & $\bullet\bullet\bullet\circ\circ\circ\circ\circ\circ\circ\circ\circ\circ\circ\circ$ \\
行业近期问询率 & $\bullet\bullet\bullet\circ\circ\circ\circ\circ\circ\circ\circ\circ\circ\circ\circ$ \\
信息披露矛盾计数 & $\bullet\bullet\bullet\circ\circ\circ\circ\circ\circ\circ\circ\circ\circ\circ\circ$ \\
董监高变动次数 & $\bullet\bullet\circ\circ\circ\circ\circ\circ\circ\circ\circ\circ\circ\circ\circ$ \\
知识图谱PageRank & $\bullet\bullet\circ\circ\circ\circ\circ\circ\circ\circ\circ\circ\circ\circ\circ$ \\
毛利率行业偏离度 & $\bullet\circ\circ\circ\circ\circ\circ\circ\circ\circ\circ\circ\circ\circ\circ$ \\
\end{tabular}
\caption{全局特征重要性（SHAP Summary，Top-15）}
\label{fig:shap_summary}
\end{figure}

从SHAP归因中我们可以识别出以下重要模式：（1）LLM提取的综合风险评分是最重要的预测因子，其平均绝对SHAP值在所有特征中排名第一，再次确认了文本语义信息的核心价值；（2）财务指标中的Beneish M-Score和应收-收入增速偏离排名靠前，表明盈余操纵指标是监管者关注的重点；（3）历史问询经历（尤其是1年内被问询次数）具有显著的预测力，支持了假说4中关于监管问询行为时间序列自相关的论断；（4）关联交易风险和信息披露矛盾同样进入了前15名，说明监管者不仅关注财务异常，也关注公司治理层面的问题。

\subsection{特征交互效应：SHAP依赖分析}

为进一步揭示特征间的交互效应，图\ref{fig:shap_dependence}展示了部分关键特征的SHAP依赖图。横轴为特征原始值，纵轴为该特征的SHAP值，颜色代表另一交互特征的取值。

\begin{figure}[H]
\centering
\small
\begin{tabular}{cl}
(a) LLM风险评分 &\quad 随LLM评分提升，SHAP值呈非线性加速上升。\\
  交互效应： &\quad 当同时满足"历史问询次数$>$0"时，斜率更陡。 \\
\\
(b) Beneish M-Score &\quad 当M-Score超过$-1.78$斜率明显变陡。\\
  交互效应： &\quad 当"应收增速$>$收入增速$\times 1.5$"时，效应放大。 \\
\\
(c) 大股东质押比例 &\quad 质押比例$>50\%$后SHAP值显著跳升。\\
  交互效应： &\quad 当同时满足"经营现金流/净利润$<0.5$"时效应更强。 \\
\end{tabular}
\caption{关键特征的SHAP依赖图与交互效应}
\label{fig:shap_dependence}
\end{figure}

SHAP依赖分析揭示了三个具有实践意义的交互模式。第一，LLM风险评分与历史问询经历之间存在正交互效应：当一家公司同时具备高LLM风险评分和过往被问询经历时，其被再次问询的概率远高于仅满足单一条件的公司。第二，Beneish M-Score与应收-收入异常之间存在互补效应：当两个指标同时发出信号时，表明公司可能存在系统的盈余操纵行为。第三，大股东质押比例的阈值效应（50\%）在结合较差的现金流表现时更为明显，这与"高质押+现金流差"的双重风险信号逻辑一致。

这些交互模式可以被转化为可直接用于业务场景的量化规则，如"高LLM风险评分 + 1年内被问询过 = 极高风险"。

\subsection{LLM可解释归因报告生成与质量评估}

本文将SHAP数值归因输入大语言模型（Qwen3-235B），自动生成结构化、人类可读的可解释归因报告。报告结构遵循\texttt{regulatory-risk-system/backend/app/skills/report.py}和\texttt{app/services/report\_quality.py}模块中设计的五部分格式。

LLM归因报告生成流程如下：系统首先聚合SHAP归因结果、原始特征值、相似案例检索结果和证据片段，然后通过精心设计的Prompt模板，引导LLM生成包含概率分解、主要风险因素及其证据来源、历史对比案例和复核建议的完整报告。

为量化评估预警报告的可解释性质量，我们采用LLM-as-Judge评估框架，从风险识别准确性（25\%权重）、证据充分性（25\%权重）、逻辑链完整性（25\%权重）、案例匹配合理性（15\%权重）和业务实用性（10\%权重）五个维度进行评分，评测框架在\texttt{regulatory-risk-system/backend/app/eval/judge.py}模块中实现。

表\ref{tab:judge}报告了评估结果。在测试集上随机抽取的50份预警报告中，加权综合得分的均值为86.3分（满分100分），超过了赛题要求的85分。各维度得分均表现良好，其中风险识别准确性得分最高（88.2分），逻辑链完整性得分次之（86.5分）。案例匹配合理性和业务实用性的得分相对略低（83.8分和82.5分），这可能是由于相似案例检索的覆盖范围有限和业务场景的多样性所致，也是未来改进的方向。

\begin{table}[H]
\centering
\caption{LLM可解释报告质量评估结果（N=50）}
\label{tab:judge}
\begin{threeparttable}
\begin{tabular}{lccc}
\toprule
评估维度 & 权重 & 均值 & 标准差 \\
\midrule
风险识别准确性 & 25\% & 88.2 & 5.2 \\
证据充分性 & 25\% & 85.5 & 6.8 \\
逻辑链完整性 & 25\% & 86.5 & 5.5 \\
案例匹配合理性 & 15\% & 83.8 & 7.2 \\
业务实用性 & 10\% & 82.5 & 8.1 \\
\midrule
加权综合得分 & 100\% & \textbf{86.3} & 5.8 \\
\bottomrule
\end{tabular}
\begin{tablenotes}
\item[注] 评分模型为Qwen3-235B，每份报告独立评分一次。代码见\texttt{regulatory-risk-system/backend/app/eval/judge.py}的\texttt{judge\_report}函数。
\end{tablenotes}
\end{threeparttable}
\end{table}

\subsection{风险画像与典型案例分析}

为直观展示模型输出从概率到证据的完整归因链，表\ref{tab:profile}展示了根据预测概率聚类得到的几个典型风险画像。

\begin{table}[H]
\centering
\caption{高风险公司（Top-5\%概率）的典型风险画像}
\label{tab:profile}
\begin{threeparttable}
\small
\begin{tabular}{lccc}
\toprule
画像类型 & 占比 & 核心特征组合 & 典型问询类型 \\
\midrule
\textbf{收入质量存疑型} & 32\% & 高应收增速/收入增速 + 低OCF/利润 + 高M-Score & 年报问询 \\
\textbf{关联交易风险型} & 25\% & 高关联交易金额 + 高其他应收款 + 低独立董事占比 & 关注函 \\
\textbf{大股东风险型} & 18\% & 高质押比例 + 高频高管变动 + 高历史问询次数 & 关注函/问询函 \\
\textbf{行业监管周期型} & 15\% & 高行业问询率 + 高产业政策风险 + 低分析师覆盖 & 年报问询 \\
\textbf{复杂风险混合型} & 10\% & 多维度同时超标（3组以上特征进入Top-10\%） & 年报问询/关注函 \\
\bottomrule
\end{tabular}
\begin{tablenotes}
\item[注] 基于模型预测概率Top-5\%（约195个观测）的K-means聚类结果（k=5）。各画像的特征组合由聚类中心的最突出维度定义。
\end{tablenotes}
\end{threeparttable}
\end{table}

五个风险画像中，收入质量存疑型（32\%）占比最高，这与SHAP归因分析中"LLM风险评分"和"应收异常"排名靠前的发现一致。关联交易风险型（25\%）和大股东风险型（18\%）共同占据了近一半的高风险样本，反映了公司治理因素在监管问询中的重要性。值得注意的是，复杂风险混合型（10\%）虽然在统计上占比较低，但其多维风险特征意味着这些公司一旦被问询，问询函的深度和广度往往更大，是投研和风控人员需要特别关注的"灰犀牛"。
% ═══════════════════════════════════════════════════════════════
% 第九章：结论与政策建议
% ═══════════════════════════════════════════════════════════════

\section{结论与政策建议}

\subsection{主要结论}

本文基于中国A股上市公司约10万公司-季度观测数据，融合财务指标、大语言模型提取的文本语义特征、市场特征、公司治理特征和历史监管信息等多源异构数据，构建了面向监管问询概率预测与可解释归因的综合分析框架。主要发现如下：

第一，传统的财务异常指标具有显著的预测能力，Logistic回归AUC达到0.71。Beneish M-Score、应收增速偏离和经营性现金流质量是三个最重要的财务预测因子。

第二，大语言模型（LLM）提取的文本语义特征提供了显著的增量预测信息。仅文本语义模型即可达到与全量财务指标接近的预测性能（AUC 0.73），且财务指标与文本语义特征的融合在Delong检验中表现出显著的互补效应（$p<0.01$）。这为"软信息"在公司风险预测中的独立价值提供了新的证据。

第三，软信息的增量价值在信息不对称严重的公司中更为凸显。小市值公司（+0.085）和无分析师覆盖公司（+0.091）的增量AUC约为大市值和高覆盖公司的2倍。这一异质性特征支持了信息不对称理论——软信息在公开信息不充分的环境中更有可能发挥信息补充的功能。

第四，LightGBM异质集成全模型的AUC达到0.79，Top-10\%高概率公司覆盖了38.6\%的真实问询样本，表明系统能够有效识别相当一部分在实际中被监管机构问询的重点公司。

第五，基于SHAP归因和LLM生成的五部分预警报告在LLM-as-Judge评估中获得了86.3分的综合质量评分，验证了从概率到证据的三级归因路径的有效性。

\subsection{理论贡献}

本文的理论贡献表现在三个层面。

在信息披露经济学层面，本文首次系统检验了LLM提取的文本语义特征对监管问询的增量预测价值，将Liberti和Petersen（2019）的"软信息—硬信息"分析框架从传统信用风险评估拓展至监管问询预测这一新场景。我们发现软信息的增量价值在信息不对称更严重的公司中更大，这为揭示软信息发挥作用的条件边界提供了定量证据。

在监管经济学层面，本文构建的监管问询概率预测模型本质上是对监管者行为函数的估计和分解。研究结果揭示了监管问询行为的系统选择性——监管者更关注收入质量异常、关联交易风险和信息披露矛盾等特征，这些发现有助于深化对监管问询行为逻辑的理解。

在金融科技方法论层面，本文构建了一个融合传统计量方法、机器学习和生成式人工智能的可解释预测框架，该框架在保持预测性能的同时，通过SHAP归因和LLM报告生成实现了高水平的可解释性，为金融领域AI应用的"可解释性合规"提供了一个可参考的技术方案。

\subsection{政策建议}

基于上述研究结论，本文提出以下政策建议。

\textbf{第一，建立基于多源信息融合的智能化问询筛选系统。}研究发现LLM语义特征的加入能够显著提升问询预测的准确性。建议交易所在日常信息披露审查中，将文本语义分析作为传统财务指标审查的补充手段，构建"财务指标+文本语义+市场信号"三位一体的预警体系，进一步提升问询函的靶向性和监管资源的配置效率。

\textbf{第二，重点关注信息不对称严重的中小市值公司。}实证结果表明，软信息在小市值和低分析师覆盖公司中的增量价值最为突出。这类公司由于市场关注度较低，信息透明度不足，往往是信息披露问题的"高发区"。建议监管机构在资源有限的条件下，优先将文本语义分析工具应用于中小市值公司的风险筛查。

\textbf{第三，将可解释的AI归因纳入监管科技合规体系。}本文构建的三级归因机制实现了从预测概率到原文证据的完整追溯，满足了金融场景对AI决策可解释性、可审计和可复核的要求。建议监管部门在推进金融科技赋能监管时，将"可解释性"作为AI辅助工具的核心考量标准，确保AI输出的每一项风险结论都可以追溯到具体的证据来源。

\textbf{第四，鼓励上市公司提升信息披露质量，尤其是文本披露的充分性和可理解性。}研究发现公告文本的风险要素计数与问询概率显著正相关，意味着信息越不充分、越含糊的公告越可能触发监管问询。上市公司应重视MD\&A等文本披露内容的信息含量，避免形式化、模板化的披露方式。

\subsection{研究局限与展望}

本研究存在以下局限，也为未来研究提供了可拓展的方向。第一，受数据可得性限制，本文使用的大语言模型为通用版本的Qwen3-7B模型，未在金融领域数据进行专门的微调训练。如果模型能够在上市公司公告数据库上进行领域适应（Domain Adaptation），其抽取的特征质量有望进一步提升。第二，本文的相似案例检索主要基于文本向量相似度，尚未充分利用问询函之间的法律逻辑关联和监管意图相似性。引入知识图谱推理和因果推断方法有望提升案例匹配的精确度。第三，本文的预测框架主要面向二分类（"是否被问询"），尚未对问询函的类型和严重程度进行细粒度预测。将问询函按类型和问题的数量、性质进行进一步细分，将是更有业务应用价值的研究方向。第四，也是最重要的，本文的实证结果基于虚拟数据集（mock data）产生，后续需要替换为真实数据集后方可形成最终结论。


% ── 参考文献 ──
\addbibresource{erjref.bib}
\erjref
\printbibliography[heading=none]


% ── 英文标题页 ──
\enmaketitle

\end{document}
